% Research-focused CV. Content is aligned with _data/cv.yml.
% Original resume template by Anubhav Singh (MIT License).
\documentclass[a4paper,10pt]{article}
\usepackage[empty]{fullpage}
\usepackage{titlesec}
\usepackage{enumitem}
\usepackage[hidelinks]{hyperref}
\usepackage{fancyhdr}
\setlength{\oddsidemargin}{-0.35in}
\setlength{\evensidemargin}{-0.35in}
\setlength{\textwidth}{7in}
\setlength{\topmargin}{-0.55in}
\setlength{\textheight}{10in}
\setlength{\parindent}{0pt}
\setlength{\parskip}{0pt}
\pagestyle{fancy}
\fancyhf{}
\fancyfoot[R]{\small Jungsoo Kim\quad\thepage}
\renewcommand{\headrulewidth}{0pt}
\renewcommand{\footrulewidth}{0pt}
\urlstyle{same}
\raggedright
\titleformat{\section}{\large\scshape}{}{0pt}{}[\titlerule]
\titlespacing*{\section}{0pt}{7pt}{3pt}
\setlist[itemize]{leftmargin=1.2em,topsep=2pt,itemsep=2pt,parsep=0pt}
\newcommand{\entry}[2]{\textbf{#1}\hfill #2\par}
\begin{document}
{\LARGE\bfseries Jungsoo Kim}\par
\vspace{1pt}
Ph.D. student in Data Science, Seoul National University\par
\href{mailto:mephisto0525@snu.ac.kr}{mephisto0525@snu.ac.kr}\enspace\textbar\enspace
\href{mailto:kimjs0619@naver.com}{kimjs0619@naver.com}\par
\href{https://www.linkedin.com/in/jungsooroot222ba/}{LinkedIn: jungsooroot222ba}\enspace\textbar\enspace
\href{https://github.com/mephistonovel}{GitHub: mephistonovel}
\section{Research Interests}

Uncertainty quantification and reliability in AI; conformal prediction and interpretable uncertainty; causal representation learning under unobserved confounding; reliability of AI-enhanced imagery. Applications span healthcare, policy evaluation, and evidentiary use.
\section{Education}
\entry{Seoul National University}{2024.09--Present}
Ph.D. in Data Science\par
Advisor: Sanghack Lee; Causality Lab.\par
\vspace{1pt}
\entry{Seoul National University}{2022.03--2024.08}
M.S. in Data Science; GPA: 4.22/4.5\par
\textbf{Course Project:} Implemented and parallelized a GRU-based name generation model from scratch using C/C++ and CUDA as part of a high-performance computing course.\par
\vspace{1pt}
\entry{Yonsei University}{2017.03--2022.02}
B.P.A. / B.B. in Science, Technology and Policy; Quantitative Risk Management\par
GPA: 4.15/4.5\par
\vspace{1pt}
\entry{National University of Singapore, Yale-NUS College}{2019.09--2020.02}
Exchange student, Mathematical and Computational Science\par
\vspace{1pt}
\section{Selected Publications}
{\small\textit{*Equal contribution (co-first authors).}}\par
\textbf{Decomposing Conformal Uncertainty: Calibration- and Instance-Driven Feature Attribution}\par
Sangyeon Cho*, Minyoung Cho*, \textbf{Jungsoo Kim}*, Sujeong Oh, Sanghack Lee.\par
\textit{NeurIPS 2026}. \textbf{Accepted, Poster}.\par
\vspace{3pt}
\textbf{A Tale of Two Uncertainties: Global-Local Attribution for Conformal Prediction}\par
Sangyeon Cho*, Minyoung Cho*, \textbf{Jungsoo Kim}*, Sujeong Oh, Sanghack Lee.\par
\textit{ICML 2026 EIML Workshop}. \textbf{Spotlight}. Earlier workshop version of the NeurIPS 2026 paper.\par
\vspace{3pt}
\textbf{Instrumental Variable Representation Learning under Confounded Covariates}\par
\textbf{Jungsoo Kim}, Kwonho Kim, Inwoo Hwang, Sanghack Lee.\par
\textit{NeurIPS 2025 CauScien Workshop}.\par
\vspace{3pt}
\textbf{When Mass Atrocities Come to Light in the Artificial Eyes in the Sky: AI and the Use of Satellite Images for International Criminal Accountability}\par
\textbf{Jungsoo Kim}, Heejin Kim, Sanghack Lee.\par
\textit{Humanities and Social Sciences Communications (Springer Nature)}. Revised and resubmitted; under review.\par
\vspace{3pt}
\section{Selected Research Projects}
\textbf{Conformal uncertainty and patient-specific retesting}\par
\begin{itemize}
\item Originated the motivating use case and led experimental design: constructed semi-synthetic data from real sepsis data and designed a risk score with patient-specific retesting needs.
\item Compared retest-variable selection using total attribution, equal calibration/instance weighting, and greater instance weighting; showed that emphasizing the test instance better prioritized patient-specific retest variables under a limited testing budget in simulation.
\end{itemize}
\vspace{1pt}
\textbf{Instrumental variable representation learning under confounded covariates}\par
\begin{itemize}
\item Conceived the research question and designed an architecture to learn IV representations from high-dimensional covariates correlated with unobserved confounders, targeting causal effect estimation under weaker assumptions.
\item Implemented the model, ran the experiments, and wrote the manuscript, carrying the project from initial concept through empirical evaluation.
\end{itemize}
\vspace{1pt}
\textbf{Reliability of AI-enhanced satellite imagery as evidence}\par
\begin{itemize}
\item Investigated the reliability of AI-enhanced satellite imagery for international criminal accountability; ran model inference across 15 super-resolution setups and conducted quantitative and qualitative analyses.
\item Combined quantitative and qualitative evidence of reliability limitations with analysis of international criminal accountability, highlighting risks of the evaluated AI-enhanced imagery as evidence; manuscript under review.
\end{itemize}
\vspace{1pt}
\newpage
\section{Research and Industrial Experience}
\entry{Ministry of Food and Drug Safety (MFDS) Project}{2025.03--2025.12}
Project Lead for Causal Inference Section\par
\begin{itemize}
\item Led the causal inference section of the MFDS project, applying conformal prediction to analyses of statin effects on blood glucose.
\item Investigated variation in prediction interval widths across individuals; this practical problem motivated subsequent research on feature attribution for conformal uncertainty.
\end{itemize}
\vspace{3pt}
\entry{SK Center for Social Value Enhancement Studies}{2021.08--2022.01}
Research Assistant\par
\begin{itemize}
\item Analyzed beneficiary financial records totaling KRW 10 billion, estimated the following fiscal year's budget, and built a research database on the Korean social economy.
\end{itemize}
\vspace{3pt}
\entry{Mooncademy}{2020.10--2021.04}
Product Manager\par
\begin{itemize}
\item Planned website information architecture with engineers and designers; launched a gayageum lecture crowdfunding campaign raising KRW 10,491,700.
\end{itemize}
\vspace{3pt}
\section{Teaching and Service}
\entry{Teaching Assistant, Seoul National University}{}
Causal Inference; Machine Learning and Deep Learning II; Basic Statistics \& Mathematics Bootcamp.\par
\vspace{5pt}
\entry{Python Instructor, SNU Language Education Institute Outreach}{2025.07--2025.08}
Taught Python and basic algorithms to Malaysian high school students, emphasizing collaborative problem-solving.\par
\vspace{5pt}
\section{Selected Honors and Writing}
\begin{itemize}
\item \textbf{5th POSTECH SF Award, Excellence Prize} (2025). Authored "Will of Stochastic Ghost," published in Diagonal Argument (Eunhengnamu, 2025); fiction exploring generative AI and evidence.
\item \textbf{Journal of Korea Artificial Intelligence Association, Best Paper} (2024). Submitted a master's project paper on instrumental variables.
\item \textbf{CSES Social Value Research Idea Award, Excellence Prize} (2024). Proposed a study on causal structure analysis for Korean population policy.
\item \textbf{LH COMPAS Land \& City Data Competition, Excellence Prize} (2024). Identified medically underserved areas in Asan through geospatial analysis.
\item \textbf{EdTech Data Analysis Hackathon, Grand Prize} (2023). Built a dropout-prediction model using survival analysis.
\item \textbf{Public BigData Analysis Hackathon, Excellence Prize} (2020). Developed a traffic-accident risk model with XGBoost and interpreted drivers using SHAP.
\end{itemize}
\section{Skills and Languages}
\textbf{Programming:} Python, R, C++, C, SQL; CUDA C/C++ (coursework)\par
\vspace{5pt}
\textbf{Research Methods:} Conformal prediction and uncertainty quantification; instrumental variables and causal representation learning; experimental design; quantitative and qualitative model evaluation\par
\vspace{5pt}
\textbf{ML and Geospatial Tools:} PyTorch; super-resolution and diffusion model workflows (SRCNN, SRGAN, StableSR, Stable Diffusion); GeoPandas, QGIS\par
\vspace{5pt}
\textbf{Data Systems:} Neo4j, BigQuery; Kafka (event streaming)\par
\vspace{5pt}
\textbf{Languages:} Korean, English (TOEFL 100; OPIc AL)\par
\vspace{5pt}
\end{document}
